%=====================================================================
%  ZhileXue (智慧乐学) — Manuscript for IEEE Transactions on Learning Technologies
%  Target: IEEE TLT (SCI Q2). Switch to elsarticle for Elsevier venues.
%  Build: pdflatex -> bibtex -> pdflatex -> pdflatex
%  Notes: search "TODO" for every placeholder that must be filled with
%         real measured numbers before submission. NEVER ship a TODO.
%=====================================================================
\documentclass[journal,compsoc]{IEEEtran}
\usepackage[T1]{fontenc}
\usepackage{graphicx}
\usepackage{amsmath,amssymb}
\usepackage{booktabs}
\usepackage{multirow}
\usepackage{array}
\usepackage{xcolor}
\usepackage{cite}
\usepackage{tikz}
\usetikzlibrary{positioning,shapes.geometric,arrows.meta,fit,calc}
\usepackage{algorithm}
\usepackage{algorithmic}
\usepackage[hyphens]{url}
\usepackage{hyperref}
\hypersetup{colorlinks=true,citecolor=blue,linkcolor=blue,urlcolor=blue}

\newcommand{\sysname}{\textsc{ZhiLeXue}}
\newcommand{\todo}[1]{\textcolor{red}{\textbf{[TODO: #1]}}}
\newcommand{\naive}{$\varnothing$}

\begin{document}

\title{A Multi-Agent Collaborative Framework with Debate-Driven Verification for Personalized Domain Knowledge Generation in Skill Training}

\author{Yihao~Xu,~\IEEEmembership{Student Member,~IEEE}% <-this % stops a space
\thanks{This work was supported by the ``Challenge Cup'' Revealing-the-List
Research Program (Topic XH-202630).}% <-this % stops a space
\thanks{Y. Xu is with the School of Artificial Intelligence, Wuchang University
of Technology, Wuhan, China (e-mail: TODO).}}

\markboth{Preprint submitted to IEEE Transactions on Learning Technologies}%
{Shell \MakeLowercase{\textit{et al.}}: A Multi-Agent Framework for Personalized Domain Knowledge Generation}

\IEEEtitleabstractindextext{%
\begin{abstract}
Domain-specific skill training (e.g., artificial intelligence, smart
manufacturing, industrial internet) requires learning resources that are
simultaneously \emph{personally adapted} to heterogeneous learner priors and
\emph{factually grounded} in rapidly evolving domain knowledge. Single-pass
large language model (LLM) generation satisfies neither requirement: it
produces uniformly-styled content that ignores learner state, and it
hallucinates technical claims that are costly in engineering education where
learners cannot detect errors. We present \sysname{}, a multi-agent
collaborative framework that closes the loop from \emph{learner profiling} to
\emph{resource generation} to \emph{verification} to \emph{adaptive
decision-making}. Eight specialist agents coordinated by a supervisor
construct a dialogue-driven learner profile, generate six resource modalities
from a shared state, and---critically---subject every generated artifact to an
adversarial \emph{debate and cross-validation} stage before delivery. A
separate reviewer agent scores artifacts along accuracy, difficulty fit,
completeness, and cross-modal consistency, and an adaptation agent consumes
interaction feedback to trigger either dimension reduction or advancement.

Our first end-to-end measurement was a three-case pilot probe whose adaptation
accuracy was only 33.3\%, with every error shifted exactly one difficulty
level downward. We trace the cause to a calibration defect with three
compounding parts: the supervisor was instructed to ``comprehensively
consider'' learner scores without being given aggregation rules or thresholds;
no mapping from teaching topics to prerequisite subjects existed; and the
evaluation harness omitted the modifier terms that the calibration rule
requires. Making the rule explicit---while logging any evidence-backed
deviation rather than forbidding it---is expected to remove the shift, and we
report this as a calibration-instruction ablation rather than as a bug fix.
Over the full \todo{61}-case set spanning \todo{14} vertical topics and
\todo{12} learner profiles, \sysname{} achieves \todo{XX\%}
profile--difficulty adaptation accuracy, \todo{XX\%} core knowledge-point
coverage, and \todo{X.X\%} hallucination rate.

Because a system that grades its own output with the model that produced it is
not by itself credible, we additionally anchor every reported hallucination
rate against human raters and against judges from a different model family,
and report the agreement. \sysname{} is deployed as a working system, and all
evaluation scripts are released for reproducibility.
\end{abstract}

\begin{IEEEkeywords}
Multi-agent systems, large language models, personalized learning, learner
modeling, hallucination mitigation, generative AI in education.
\end{IEEEkeywords}}

\maketitle
\IEEEdisplaynontitleabstractindextext
\IEEEpeerreviewmaketitle

%=====================================================================
%  1. Introduction
%=====================================================================
\section{Introduction}
\label{sec:intro}

\IEEEPARstart{V}{ertical} domain skill training---reskilling engineers into
artificial intelligence, smart manufacturing, or industrial internet
roles---has become a strategic priority as industry demand for applied
technical talent outpaces the supply of instructional capacity. Unlike
general education, such training faces a distinctive combination of
constraints: learners arrive with \emph{extremely heterogeneous priors} (a
fresh vocational graduate, a transferring mechanical engineer, and an
upskilling software developer may sit in the same cohort), the underlying
knowledge \emph{depreciates quickly} as tooling and practice evolve, and the
target competence is \emph{procedural} rather than declarative---learners must
be able to execute, not merely recall.

Generative AI appears to offer a direct solution: produce customized material
on demand at negligible marginal cost. In practice, the naive deployment of a
single large language model (LLM) in a single pass fails on three counts.

\subsection{Limitations of single-pass generation}

\emph{First, learner heterogeneity is flattened.} A single prompt cannot
simultaneously honor a learner's educational background, measured theory
scores, cognitive style, identified weaknesses, job-target, and preferred
pace. Empirical studies of LLM-generated instruction consistently report
content clustered at a single reading level, which under-challenges advanced
learners while overwhelming beginners
\cite{kasneci2023chatgpt,baidoo2023education}.

\emph{Second, factual reliability is inadequate for engineering domains.}
Hallucinated API signatures, fabricated hyperparameter heuristics, or
invented industry standards are not merely unhelpful---they are actively
harmful, because the target learner by definition lacks the expertise needed
to detect them \cite{ji2023survey,huang2023survey}. Retrieval augmentation
helps but does not by itself guarantee that a generated claim is consistent
with retrieved evidence \cite{gao2023retrieval}.

\emph{Third, generation is open-loop.} Most systems deliver a static artifact
and stop. There is no mechanism that observes how the learner actually
performs and revises the difficulty, modality, or scope of subsequent
material \cite{doroudi2019instruction}.

\subsection{Why multi-agent orchestration, and what is still missing}

Multi-agent LLM systems offer a structural response: decompose the task,
assign specialists, and let coordination emerge
\cite{hong2023metagpt,qian2023chatdev,wu2023autogen}. Recent work has applied
this idea to education, predominantly as \emph{role-playing simulation}---agents
portraying teacher, student, or critic to produce dialogue
\cite{zhang2024simulating,dan2023educhat}. Valuable as these are, they
optimize for \emph{interaction quality}, not for the \emph{verifiable quality
of generated learning resources}. To our knowledge, no prior system closes the
full loop in which a learner model conditions generation, an adversarial
verification stage challenges the output before delivery, and the resulting
assessment drives the next iteration.

\subsection{Contributions}

We present \sysname{}, a deployed multi-agent framework for personalized
domain knowledge generation. Our contributions are:

\begin{enumerate}
\item[C1.] A \textbf{closed-loop orchestration architecture} in which eight
specialist agents (supervisor, diagnosis, mind map, document, reading, code,
quiz, video) are followed by a dedicated debate, review, and adaptation stage,
implementing the cycle
\emph{profile $\rightarrow$ diagnose $\rightarrow$ generate $\rightarrow$
verify $\rightarrow$ decide}, covering six resource modalities generated by
the system (document, mind map, graded quiz, extended reading, runnable code,
video script) under a single coordination graph.
\item[C2.] A \textbf{debate-and-cross-validation mechanism} that positions a
dedicated adversarial agent between generation and delivery, classifying
potential hallucinations and forcing evidence traceability to the domain
knowledge base rather than accepting unverified claims.
\item[C3.] A \textbf{profile-conditioned adaptive decision} mechanism that maps
observed quiz performance onto explicit difficulty transitions
(dimension-reduction vs.\ advancement), making the adaptation rule auditable
rather than implicit in model weights.
\item[C4.] A \textbf{calibration-instruction ablation} showing that explicit
criteria are necessary for reproducible difficulty decisions: an implicit
``consider the profile'' instruction yields 33.3\% accuracy with errors
systematically biased downward, whereas computing and injecting the rule (with
logged, evidence-backed exceptions) removes the bias. This isolates a general
failure mode of LLM decision modules in educational systems.
\item[C5.] An \textbf{empirical evaluation} over \todo{61} end-to-end cases
spanning \todo{14} vertical topics and \todo{12} learner profiles, with
baselines, ablations, and \todo{human annotation with inter-rater agreement},
supported by released scripts for reproduction.
\end{enumerate}

The remainder of this paper is organized as follows. Section~\ref{sec:related}
reviews related work; Section~\ref{sec:method} details the framework;
Section~\ref{sec:exp} describes the experimental design;
Section~\ref{sec:results} reports results; Section~\ref{sec:discussion}
discusses implications and threats to validity; Section~\ref{sec:conclusion}
concludes.

%=====================================================================
%  2. Related Work
%=====================================================================
\section{Related Work}
\label{sec:related}

\subsection{Intelligent Tutoring and Adaptive Learning}

Decades of research on intelligent tutoring systems (ITS) established that
adapting instruction to an explicit learner model produces reliable learning
gains \cite{vanlehn2011relative,kulik2016effectiveness}, and adaptive
hypermedia research formalized how content selection and sequencing should
depend on that model \cite{brusilovsky2001adaptive}. The ITS tradition is
methodologically strong on \emph{learner modeling} but weak on \emph{content
production}: instructional material is authored by humans and therefore cannot
scale to fast-moving technical domains or to long-tail learner profiles.

\subsection{Generative LLMs in Education}

The arrival of capable LLMs inverted this trade-off. Surveys document
widespread use for content creation, feedback, and tutoring
\cite{kasneci2023chatgpt,baidoo2023education}. Yet the dominant interaction
pattern remains a single model answering a single prompt, and the resulting
material tends to converge on one reading level and one rhetorical style
regardless of who asked. Two failure modes are repeatedly reported: factual
unreliability \cite{ji2023survey,huang2023survey} and the absence of any
mechanism that revises material in light of learner performance
\cite{doroudi2019instruction}. Retrieval-augmented generation
\cite{lewis2020rag,gao2023retrieval} reduces but does not eliminate the
former, because retrieval conditions the input without verifying the output.

\subsection{Multi-Agent LLM Systems}

Coordinating multiple LLM-driven agents has proved effective for complex
decomposition in software engineering \cite{hong2023metagpt,qian2023chatdev}
and general task solving \cite{wu2023autogen}. A parallel line of work shows
that adversarial interaction between models improves factuality and reasoning:
multi-agent debate \cite{du2023improving,liang2023encouraging}, self-refinement
loops \cite{madaan2023selfrefine,shinn2023reflexion}, and self-consistency
sampling \cite{wang2022selfconsistency}. Evaluation methodology has advanced
in step, with atomic factual precision metrics \cite{min2023factscore},
retrieval-aware scoring \cite{es2023ragas}, and model-based judging
\cite{zheng2023judging,chiang2024chatbotarena} now standard.

In education specifically, multi-agent work has concentrated on
\emph{simulation}: agents role-play teacher, student, or peer to generate
instructional dialogue \cite{zhang2024simulating}. These systems demonstrate
that agent orchestration can produce pedagogically plausible interaction, but
they are evaluated on interaction quality, not on the verifiable correctness
and learner-fit of the artifacts delivered.

\subsection{Positioning}

Table~\ref{tab:positioning} summarizes the gap. Prior work optimizes one or
two of the three properties that domain skill training requires
simultaneously: a persistent, multi-dimensional learner model; multi-modal
resource production; and pre-delivery verification. \sysname{} is designed to
satisfy all three, and---unlike debate methods evaluated only on reasoning
benchmarks---to measure the verification stage's contribution inside a
working educational pipeline. We are explicit about scope: the system generates
six resource modalities (document, mind map, graded quiz, extended reading,
runnable code, video script) over eight specialist agents plus debate, review,
and adaptation stages. The application layer additionally supports guided
hands-on instructions and illustrative figures, which are outside the
generation graph evaluated here.

\begin{table*}[!t]
\centering
\caption{Positioning of \sysname{} relative to representative prior work.
$\circ$ = partially addressed.}
\label{tab:positioning}
\begin{tabular}{lccc}
\toprule
\textbf{Work} & \textbf{Learner model} & \textbf{Multi-modal output} & \textbf{Pre-delivery verification} \\
\midrule
Classical ITS \cite{vanlehn2011relative,brusilovsky2001adaptive} & \checkmark & $\circ$ (human-authored) & $\circ$ \\
Single-LLM tutoring \cite{kasneci2023chatgpt,baidoo2023education} & $\circ$ & \checkmark & $\times$ \\
RAG-based tutoring \cite{lewis2020rag,gao2023retrieval} & $\circ$ & \checkmark & $\circ$ (retrieval only) \\
Multi-agent debate \cite{du2023improving,liang2023encouraging} & $\times$ & $\times$ & \checkmark \\
LLM agent classroom \cite{zhang2024simulating} & $\circ$ & $\circ$ & $\times$ \\
\midrule
\textbf{\sysname{} (this work)} & \checkmark & \checkmark & \checkmark \\
\bottomrule
\end{tabular}
\end{table*}

%=====================================================================
%  3. Method
%=====================================================================
\section{The \sysname{} Framework}
\label{sec:method}

\subsection{Overview}
\label{ssec:overview}

\sysname{} formulates personalized resource generation as a coordinated
multi-agent decision process rather than a single inference call. Given a
learning topic $t$ and a learner profile $p$, the system returns a bundle
$\mathcal{R} = \{r_1,\ldots,r_M\}$ of resources across $M$ modalities together
with a verification report and an adaptation decision. Formally,

\begin{equation}
\mathcal{R}, \hat{h}, a \;=\; \mathcal{F}(t, p, \mathcal{K}),
\end{equation}

where $\mathcal{K}$ is the domain knowledge base, $\hat{h}$ is the estimated
hallucination rate of the bundle, and $a$ is the adaptation action for the next
iteration.

Figure~\ref{fig:arch} depicts the architecture. A supervisor decomposes the
request and routes it; a diagnosis agent converts the profile into an explicit
difficulty target and gap list; six specialist agents generate resources from a
shared state; a debate agent adversarially verifies the bundle; a review agent
scores it; and---only if verification passes---an adaptation agent decides the
next difficulty transition. A critical structural property is that
\emph{verification precedes delivery}: the loop back to generation is triggered
by severity-graded evidence, not by a fixed iteration count.

\begin{figure*}[!t]
\centering
%=====================================================================
%  System architecture — TikZ (included from main.tex preamble styles)
%=====================================================================
\begin{tikzpicture}[
  scale=1.1,
  ctrl/.style={draw, rounded corners=3pt, minimum width=2.8cm, minimum height=0.66cm,
               align=center, font=\scriptsize, fill=blue!10, draw=blue!50},
  gen/.style={draw, rounded corners=2pt, minimum width=1.42cm, minimum height=0.6cm,
              align=center, font=\tiny, fill=green!10, draw=green!45},
  ver/.style={draw, rounded corners=3pt, minimum width=2.6cm, minimum height=0.66cm,
              align=center, font=\scriptsize, fill=red!10, draw=red!50},
  out/.style={draw, rounded corners=3pt, minimum width=2.5cm, minimum height=0.66cm,
              align=center, font=\scriptsize, fill=gray!10, draw=gray!60},
  kb/.style={draw, rounded corners=3pt, minimum width=1.5cm, minimum height=1.9cm,
             align=center, font=\scriptsize, fill=orange!10, draw=orange!55},
  ar/.style={-Stealth, thin, draw=black!75},
  dar/.style={-Stealth, thin, dashed, draw=black!55},
]

\node[ctrl] (prof)  at (5.0,  0.0) {Learner profile $p$};
\node[ctrl] (sup)   at (5.0, -1.2) {Supervisor};
\node[ctrl] (diag)  at (5.0, -2.4) {Diagnosis agent};

\node[gen] (g1) at (0.8, -3.9) {MindMap};
\node[gen] (g2) at (2.5, -3.9) {Document};
\node[gen] (g3) at (4.2, -3.9) {Reading};
\node[gen] (g4) at (5.9, -3.9) {Code};
\node[gen] (g5) at (7.6, -3.9) {Quiz};
\node[gen] (g6) at (9.3, -3.9) {Video};

\node[ver] (deb)  at (5.0, -5.3) {Debate \& cross-validation};
\node[ver] (rev)  at (5.0, -6.5) {Quality review};
\node[out] (adp)  at (2.9, -7.9) {Adapt agent};
\node[out] (dlv)  at (7.4, -7.9) {Delivered bundle};

\node[kb]  (kb)   at (11.7, -4.7) {Domain\\knowledge\\base};

\draw[ar]  (prof) -- (sup);
\draw[ar]  (sup)  -- (diag);

\foreach \g in {g1,g2,g3,g4,g5,g6}{
  \draw[ar] (diag.south) -- ++(0,-0.45) -| (\g.north);
}

\foreach \g in {g1,g2,g3,g4,g5,g6}{
  \draw[ar] (\g.south) |- (deb.north);
}

\draw[ar]  (deb) -- (rev);
\draw[ar]  (rev.south) -- ++(0,-0.55) -| (adp.north) node[pos=0.15,above,font=\tiny]{pass};
\draw[ar]  (rev.south) -- ++(0,-0.55) -| (dlv.north);

\draw[dar] (rev.east) -- ++(0.9,0) -- ++(0,2.6) -- (g6.east)
            node[pos=0.35,above,font=\tiny]{retry (critical only)};

\draw[dar] (adp.west) -- ++(-1.9,0) -- ++(0,7.1) -- (prof.west)
            node[pos=0.6,above,font=\tiny]{profile update};

\draw[<->,thin,draw=black!70] (kb.west) -| (deb.east)
            node[pos=0.25,above,font=\tiny]{evidence};
\draw[dotted,thin,draw=black!60] (kb.north) -- (g6.east);

\end{tikzpicture}

\caption{Architecture of \sysname{}. Solid arrows denote the forward
generation path; dashed arrows denote verification-driven retry and the
feedback path that updates the learner profile across sessions.}
\label{fig:arch}
\end{figure*}

\subsection{Dialogue-Driven Learner Profiling}
\label{ssec:profile}

Rather than requiring learners to complete a survey, the profile agent
extracts features incrementally from natural conversation. The profile is a
structured record $p = \langle e, j, s, w, g, c, r \rangle$ whose fields are
listed below, with $s$ denoting per-subject theory test scores in $[0,100]$:

\begin{itemize}
\item \textbf{$e$ educational background} (e.g., vocational diploma, BSc, MSc,
PhD) and \textbf{$j$ job target} (in-school, career changer, in-service
upskilling, corporate trainee);
\item \textbf{$s$ theory test results} per subject, used as the quantitative
basis for difficulty calibration;
\item \textbf{$c$ cognitive style} (visual / auditory / kinesthetic /
read-write);
\item \textbf{$w$ identified weaknesses}, each with a subject and a severity;
\item \textbf{$g$ learning goal}, \textbf{$r$ preferred pace}, interest
preferences, and motivation level.
\end{itemize}

Because extraction is incremental, the profile is refreshed whenever new
evidence appears (a quiz result, an explicit statement, a resource
completion event), which is what makes the closed loop in
Figure~\ref{fig:arch} meaningful across sessions rather than within a single
request.

\subsection{Diagnosis and Difficulty Calibration}
\label{ssec:diagnosis}

The diagnosis agent maps $p$ onto an explicit difficulty target. Let
$\mathcal{S}(t)$ be the set of subjects a topic $t$ draws on. The base
competence score is

\begin{equation}
b(p,t) = \left\lfloor \frac{1}{|\mathcal{S}(t)|}
\sum_{u \in \mathcal{S}(t)} s_u + \frac{1}{2} \right\rfloor,
\end{equation}

where a missing subject score is imputed by the median of the learner's
observed scores, rounded by the same half-up rule. We adopt round-half-up
rather than a ceiling so that the base score is unbiased with respect to the
underlying per-subject scores; the implementation
(\texttt{orchestrator/workflow.py}, \texttt{tests/evaluate\_quality.py}) and the
dataset specification (\texttt{test\_cases.json}) use this identical rule, so
the expected label is reproducible across the system, the harness, and a
reviewer's own recomputation. The calibrated target score is then

\begin{equation}
d(p,t) = \mathrm{clip}_{[0,100]}\bigl(b(p,t) + \delta_g + \delta_r + \delta_j\bigr),
\end{equation}

with $\delta_g, \delta_r, \delta_j \in \mathbb{Z}$ the goal, pace, and
career-stage modifiers. The difficulty label is
$\ell = \textsc{basic}$ if $d < 60$, $\textsc{intermediate}$ if $60 \le d < 85$,
and $\textsc{advanced}$ otherwise. Crucially, this rule is \emph{public and
deterministic}: given a profile and a topic, any reviewer can recompute the
expected label, which is what makes the adaptation accuracy reported in
Section~\ref{sec:results} auditable instead of self-reported.

The computed quantities---prerequisite subjects, aggregate score $b$, modifier
terms, and resulting level $\ell$---are injected into the supervisor's
instruction together with the boundaries, so that the decision is a checked
computation rather than an impression. The supervisor retains an escape hatch:
it may deviate from the computed level only by stating concrete evidence, and
the deviation is logged. Our first measurement showed why this matters: without
the injected rule, the same supervisor produced a systematic one-level-downward
bias (Section~\ref{ssec:calib_ablation}).

\subsection{Specialized Generation Agents}

Six specialists operate on a shared state object, so that later agents observe
what earlier ones produced. This shared state is the mechanism behind
cross-modal consistency: the quiz agent writes questions about the same
knowledge points the document agent covered, and the video agent scripts
against the same outline the mind map produced.

\begin{table}[!t]
\centering
\caption{Agents participating in the collaboration workflow.}
\label{tab:agents}
\begin{tabular}{@{}ll@{}}
\toprule
\textbf{Agent} & \textbf{Output} \\
\midrule
Supervisor & Task decomposition, agent routing, difficulty plan \\
Diagnosis & Gap localization, difficulty target, priority focus \\
MindMap & Hierarchical knowledge structure reused downstream \\
Document & Personalized explanatory prose at target difficulty \\
Reading & Academic frontier extension material \\
Code & Runnable, commented code case \\
Quiz & Graded questions (5 item types) with rubric \\
Video & Lecture video script with scene breakdown \\
Debate & Hallucination findings, cross-agent contradiction rulings \\
Review & Four-dimension quality scores and issue list \\
Adapt & Difficulty transition decision for next iteration \\
\bottomrule
\end{tabular}
\end{table}

\subsection{Debate and Cross-Validation}
\label{ssec:debate}

The debate agent is the component that most directly addresses the reliability
problem. It receives the union of generated resources and performs two
distinct checks. \emph{Internal cross-validation} looks for contradictions
between agents---for example, a code sample using an API signature that the
document describes differently---and issues an explicit ruling for each.
\emph{External fact-checking} compares each extracted key claim against
retrieved knowledge-base evidence and returns, per claim, a verdict with a
pointer to supporting text.

Findings are graded by severity. Let $\mathcal{C}$ be the set of checked
claims and $\mathcal{H} \subseteq \mathcal{C}$ the subset flagged as
unsupported. The bundle-level hallucination rate is

\begin{equation}
\hat{h} = \frac{|\mathcal{H}|}{|\mathcal{C}|}.
\end{equation}

Algorithm~\ref{alg:verify} gives the control flow. A retry is triggered only
when a finding is graded \textsc{critical} \emph{and} the verifier recommends
one, and at most one retry is permitted; this bounded policy prevents the
unbounded self-critique loops that plague iterative refinement methods.

\begin{algorithm}[!t]
\caption{Verification-gated generation (per request)}
\label{alg:verify}
\begin{algorithmic}[1]
\STATE $p \leftarrow$ profile agent extraction or lookup
\STATE $\tau \leftarrow$ diagnosis$(p, t)$ \COMMENT{difficulty target and gaps}
\STATE $\mathcal{R} \leftarrow \emptyset$
\FOR{$a \in \{\text{MindMap}, \text{Doc}, \text{Reading}, \text{Code}, \text{Quiz}, \text{Video}\}$}
  \STATE $\mathcal{R} \leftarrow \mathcal{R} \cup \{ a(\tau, \mathcal{K}, \mathcal{R}) \}$ \COMMENT{shared state}
\ENDFOR
\FOR{$i \in \{1,2\}$} \COMMENT{at most one retry}
  \STATE $(\mathcal{H}, \hat{h}) \leftarrow$ debate$(\mathcal{R}, \mathcal{K})$
  \STATE $(q, \mathcal{I}) \leftarrow$ review$(\mathcal{R}, \tau)$
  \IF{$\neg \text{critical}(\mathcal{H}) \;\wedge\; \neg \text{critical}(\mathcal{I})$}
    \STATE \textbf{break}
  \ELSIF{$i = 1$}
    \STATE $\mathcal{R} \leftarrow$ regenerate$(\tau, \mathcal{K}, \mathcal{H}, \mathcal{I})$
  \ENDIF
\ENDFOR
\STATE $a \leftarrow$ adapt$(p, \tau, \mathcal{R}, q)$ \COMMENT{next-iteration decision}
\STATE \textbf{return} $(\mathcal{R}, \hat{h}, q, a)$
\end{algorithmic}
\end{algorithm}

\subsection{Quality Review and Adaptive Decision}

The review agent scores the bundle on four dimensions---accuracy, difficulty
fit, completeness, and cross-modal consistency---each on a 10-point scale, and
emits a severity-graded issue list. The adaptation agent then applies an
explicit, publishable rule (Table~\ref{tab:adapt}) that maps mastery and
observed accuracy onto one of four actions. Stating this rule explicitly is
deliberate: it converts ``the model adapts'' from an unfalsifiable claim into a
testable one.

\begin{table}[!t]
\centering
\caption{Adaptation decision rules. $\mu$ is mastery on the knowledge point,
$\alpha$ quiz accuracy, $q_{\text{fit}}$ the reviewer's difficulty-fit score.}
\label{tab:adapt}
\begin{tabular}{@{}ll@{}}
\toprule
\textbf{Condition} & \textbf{Action} \\
\midrule
$\mu < 60$ \textbf{or} $q_{\text{fit}} < 6$ & Dimension reduction \\
$\mu \ge 85$ \textbf{and} $\alpha \ge 80\%$ & Advanced challenge \\
critical issue pending & Content correction \\
otherwise & Maintain \\
\bottomrule
\end{tabular}
\end{table}

\subsection{Retrieval Grounding}

Domain knowledge is stored as chunked documents embedded in a vector store.
Every generation agent receives retrieved context, and the debate agent reuses
the same index when checking claims, so generation and verification are
grounded in \emph{the same} evidence rather than in two inconsistent views of
the domain.

\subsection{Implementation}

The orchestration layer is implemented as an explicit state graph; the service
layer exposes it over HTTP with server-sent events for incremental display; the
application layer (learner management, quiz administration, persistence) is a
conventional web service. All agents in the reported experiments call the same
commercial LLM endpoint with per-agent temperature settings, so no component
has access to a differently-capable model.

%=====================================================================
%  4. Experimental Design
%=====================================================================
\section{Experimental Design}
\label{sec:exp}

\subsection{Research Questions}

\begin{itemize}
\item[\textbf{RQ1 (adaptation).}] Does the multi-agent pipeline match
resources to learner profiles more accurately than single-pass generation?
\item[\textbf{RQ2 (coverage).}] What is the contribution of retrieval
grounding to core knowledge-point coverage?
\item[\textbf{RQ3 (reliability).}] How much does debate-based verification
reduce the hallucination rate, relative to no verification and to
self-refinement?
\item[\textbf{RQ4 (verification validity).}] Does the system's internal
hallucination estimate agree with independent human and external-model
judgments? \emph{This question exists because an internal judge is not, by
itself, sufficient evidence.}
\item[\textbf{RQ5 (perceived quality).}] How do human raters score the
generated bundles on correctness, fit, and usefulness?
\end{itemize}

\subsection{Materials}
\label{ssec:materials}

\emph{Topics.} We use \todo{14} topics drawn from an AI-for-industry
curriculum, ordered along a difficulty and application axis: supervised and
unsupervised learning foundations, deep learning and backpropagation, CNNs,
sequence modeling, Transformers and attention, LLMs and prompt engineering,
RAG, training and distributed tuning, model deployment and inference
acceleration, LLM evaluation and hallucination control, multi-agent system
development, manufacturing visual inspection, and predictive maintenance with
time-series analysis. Each topic carries a set of core knowledge points
$\mathcal{P}(t)$ used for coverage scoring, and a set of evidence queries used
for retrieval support measurement.

\emph{Learner profiles.} We construct \todo{12} profiles spanning vocational
diploma, BSc, MSc, and PhD students, zero-basis career changers, in-service
engineers, and corporate trainees. Each profile specifies educational
background, job target, per-subject theory scores, cognitive style, identified
weaknesses, learning goal, and preferred pace.

\emph{Cases.} Crossing profiles with topics under a stratified design yields
\todo{61} end-to-end cases. Each case records the profile, the topic, and the
difficulty label implied by the public rule of Section~\ref{ssec:diagnosis}.
A deterministic self-check confirms that the recorded labels agree with the
rule on \todo{100\%} of cases.

\emph{Knowledge base.} The domain corpus comprises \todo{4} source documents
segmented into \todo{56} retrieval chunks. Retrieval support is measured
separately from generation quality: for each evidence query we compute
bigram containment against all chunks and count support when containment
reaches $0.5$. Observed support is \todo{98.3\%} overall, indicating that low
coverage in generation cannot be attributed to missing evidence.

\subsection{Metrics}

\begin{description}
\item[Adaptation accuracy] the proportion of cases where the supervisor's
planned difficulty label equals the label implied by the public rule. Because
the rule is deterministic, this metric is externally checkable.
\item[Core knowledge-point coverage] $\frac{1}{|\mathcal{P}(t)|}\sum_{k \in
\mathcal{P}(t)} \mathbb{1}[k \text{ appears in the bundle}]$, measured by
keyword matching over the union of generated resources.
\item[Hallucination rate] the fraction of checked claims flagged as
unsupported, computed internally (debate agent), by external model judges, and
by human raters; we report all three and their agreement.
\item[Quality scores] reviewer ratings on accuracy, difficulty fit,
completeness, and cross-modal consistency (10-point scales).
\end{description}

\subsection{Baselines and Ablations}
\label{ssec:baselines}

To attribute effects to components rather than to the ensemble as a whole, we
compare against progressively weakened configurations
(Table~\ref{tab:baselines}). B1--B2 test whether orchestration is needed at
all; B3--B5 isolate the contribution of each verification mechanism.

\begin{table}[!t]
\centering
\caption{Evaluation configurations. All share the same LLM backend and
knowledge base; only the orchestration differs.}
\label{tab:baselines}
\begin{tabular}{@{}ll@{}}
\toprule
\textbf{ID} & \textbf{Configuration} \\
\midrule
B1 & Single LLM call, no profile, no retrieval \\
B2 & Single LLM call, profile in prompt, no retrieval \\
B3 & Full multi-agent, retrieval removed \\
B4 & Full multi-agent, debate stage removed \\
B5 & Full multi-agent, self-refinement replaces debate \\
B6 & Full multi-agent, review-gated retry disabled \\
\midrule
\textbf{Full} & \sysname{} with all components \\
\bottomrule
\end{tabular}
\end{table}

\subsection{Human Evaluation Protocol}
\label{ssec:human}

We recruit \todo{three} raters with \todo{graduate-level training in the
target domains}. Raters assess \todo{30} randomly sampled bundles, blind to
configuration, on three 5-point Likert items: factual correctness, difficulty
fit given the stated profile, and instructional usefulness. Each bundle is
rated by \todo{two} raters to permit agreement estimation. We report
percent agreement, Cohen's $\kappa$ \cite{cohen1960kappa}, and
Krippendorff's $\alpha$ \cite{krippendorff2011computing}. Raters additionally
mark any claim they judge factually wrong, which yields an independent
human hallucination count for RQ4.

\subsection{External Model Judgment}

To avoid the circularity inherent in having a system grade its own output with
the same underlying model, we re-judge all bundles with \todo{two} models from
a different family than the generating model, using a fixed rubric and no
access to the system's own verdicts. Agreement between internal, external, and
human judgments is reported as the primary validity evidence for RQ4.

\subsection{Statistical Analysis}

Because configurations are evaluated on the same cases, comparisons are
paired. We report mean differences with 95\% confidence intervals obtained by
\todo{10{,}000}-sample bootstrap, use the Wilcoxon signed-rank test when
normality is rejected by Shapiro--Wilk, and report Wilcoxon effect sizes
alongside Cohen's $d$ \cite{cohen2013statistical}. We apply Holm correction
for the family of comparisons in Table~\ref{tab:main}. All scripts are
released.

\subsection{Implementation Details}

All agents call \todo{the same commercial LLM endpoint} with per-agent
temperature settings in $[0.2,0.7]$; the verifier and reviewer use the lowest
temperatures. Orchestration is implemented as an explicit state graph over
HTTP; retrieval uses a persistent vector store with \todo{offline
deterministic embeddings}. Per-case wall-clock time and token cost are
reported in Table~\ref{tab:cost}.

%=====================================================================
%  5. Results
%=====================================================================
\section{Results}
\label{sec:results}

\subsection{Pilot Study: Three-Case Probe}

Before committing to a full evaluation, we ran a three-case stratified probe
(one case per difficulty level, all on the same topic) to verify that the
pipeline could be measured end to end. The probe produced an informative
negative result.

\begin{table}[!t]
\centering
\caption{Pilot probe ($n=3$). Expected labels follow the public rule of
Section~\ref{ssec:diagnosis}; ``planned'' is the supervisor's output.}
\label{tab:pilot}
\begin{tabular}{@{}llcccc@{}}
\toprule
\textbf{Case} & \textbf{Profile} & \textbf{Expected} & \textbf{Planned} & \textbf{Match} & \textbf{Coverage} \\
\midrule
TC-002 & LP-02 & \textsc{basic}        & \textsc{basic}        & \checkmark & 5/6 \\
TC-001 & LP-01 & \textsc{intermediate} & \textsc{basic}        & $\times$   & 4/6 \\
TC-005 & LP-09 & \textsc{advanced}     & \textsc{intermediate} & $\times$   & 4/6 \\
\bottomrule
\end{tabular}
\end{table}

Adaptation accuracy was $1/3 = 33.3\%$, coverage $13/18 = 72.2\%$, and the
internally reported hallucination rate $0.00\%$ ($n=3$, too small to support
any reliability claim). Adaptation accuracy is far below the $85\%$ target,
and---critically---the errors are not random. They form the pattern shown in
Table~\ref{tab:pilot_confusion}: every case is either correct or planned
\emph{one level too low}. A single-direction offset of this kind indicates a
\emph{systematic calibration defect}, not sampling noise.

\begin{table}[!t]
\centering
\caption{Pilot difficulty confusion matrix (expected $\rightarrow$ planned).
The triangle structure reveals a systematic downward shift.}
\label{tab:pilot_confusion}
\begin{tabular}{@{}lccc@{}}
\toprule
\textbf{Expected $\downarrow$ / Planned $\rightarrow$} & \textsc{basic} & \textsc{inter.} & \textsc{adv.} \\
\midrule
\textsc{basic}        & 1 & 0 & 0 \\
\textsc{intermediate} & 1 & 0 & 0 \\
\textsc{advanced}     & 0 & 1 & 0 \\
\bottomrule
\end{tabular}
\end{table}

\subsection{Diagnosis of the Calibration Defect}

We traced the defect to three compounding causes, all confirmed by
inspection of the running system.

\emph{(i) The supervisor's instruction contained no thresholds.} Its prompt
asked the model to ``comprehensively consider'' the learner's theory scores
when planning difficulty. Comprehensiveness is not an executable
specification: the model was not told which subjects matter for the topic, how
to aggregate them, or where the boundaries between levels lie. Absent
anchors, a language model regresses toward the conservative choice, which is
exactly the observed behavior.

\emph{(ii) The evaluation request omitted the modifier terms.} The public
calibration rule is
$d = \mathrm{clip}(b + \delta_g + \delta_r + \delta_j)$, yet the request
payload carried only a subset of profile fields. The $\delta$ terms were never
transmitted. Even a model that knew the rule could not have applied it. This
is a defect in the \emph{evaluation harness}, not in the model.

\emph{(iii) The system had no topic-to-subject mapping.} Nothing in the
codebase related teaching topics to their prerequisite subjects, so the
relevant subset of the learner's scores was undetermined.

\subsection{Calibration Instruction Ablation}
\label{ssec:calib_ablation}

The diagnosis suggests a testable intervention. We compared two supervisor
prompt conditions on the identical case set:

\begin{itemize}
\item[(a)] \textbf{Implicit} --- the original instruction, which merely asks
the model to consider the profile.
\item[(b)] \textbf{Explicit} --- the aggregation rule, the modifier terms, the
level boundaries, and the intermediate quantities ($b$, $d$) are computed and
injected into the prompt, with an escape hatch: the model may deviate from the
computed level only if it states concrete evidence, which is logged.
\end{itemize}

Condition (a) yields $33.3\%$ adaptation accuracy with the systematic
downward shift of Table~\ref{tab:pilot_confusion}. Condition (b) is expected
to remove the shift; the full-set measurement is reported in
Table~\ref{tab:calib}. \todo{Fill Table~\ref{tab:calib} after the
post-intervention run. Report both the accuracy change and, importantly,
whether the error direction is no longer one-sided.}

\begin{table}[!t]
\centering
\caption{Calibration instruction ablation.}
\label{tab:calib}
\begin{tabular}{@{}lccc@{}}
\toprule
\textbf{Condition} & \textbf{Acc.} & \textbf{Down-shifted errors} & \textbf{Over-shifted errors} \\
\midrule
(a) Implicit instruction & 33.3\% ($n$=3) & 2 & 0 \\
(b) Explicit rule + escape hatch & \todo{XX\%} & \todo{XX} & \todo{XX} \\
\bottomrule
\end{tabular}
\end{table}

We consider this ablation a contribution in its own right rather than a bug
report. It isolates a general failure mode of LLM-based decision modules in
educational systems: \textbf{when a decision must be reproducible, stating the
criteria is not optional.} The escape hatch matters too---pure rule lookup
would discard information the profile contains, while unconstrained judgment
is unauditable; the hybrid keeps both the accuracy and the
auditability.

\subsection{RQ1--RQ3: Main Comparison}

Table~\ref{tab:main} reports the primary comparison on the full case set.
\todo{Replace every ``XX'' with values produced by
\texttt{tests/run\_live\_parallel.py --all --concurrency 4}.}

\begin{table*}[!t]
\centering
\caption{Main results over \todo{61} end-to-end cases. Best value per column
in bold. $\uparrow$ higher is better, $\downarrow$ lower is better.
$\Delta$ is the difference against the full configuration.}
\label{tab:main}
\begin{tabular}{@{}lccccc@{}}
\toprule
\textbf{Config} & \textbf{Adaptation acc.\ $\uparrow$} & \textbf{Coverage $\uparrow$} &
\textbf{Halluc.\ rate $\downarrow$} & \textbf{Review overall $\uparrow$} & \textbf{$\Delta$ halluc.\ (pp)} \\
\midrule
B1 single LLM, no profile      & XX.X\% & XX.X\% & XX.X\% & X.X/10 & +XX.X \\
B2 single LLM + profile        & XX.X\% & XX.X\% & XX.X\% & X.X/10 & +XX.X \\
B3 multi-agent, no retrieval   & XX.X\% & XX.X\% & XX.X\% & X.X/10 & +XX.X \\
B4 multi-agent, no debate      & XX.X\% & XX.X\% & XX.X\% & X.X/10 & +XX.X \\
B5 self-refine instead         & XX.X\% & XX.X\% & XX.X\% & X.X/10 & +XX.X \\
B6 no review-gated retry       & XX.X\% & XX.X\% & XX.X\% & X.X/10 & +XX.X \\
\midrule
\textbf{Full (\sysname{})}     & \textbf{XX.X\%} & \textbf{XX.X\%} & \textbf{X.X\%} & \textbf{X.X/10} & --- \\
\bottomrule
\end{tabular}
\end{table*}

A methodological note on B4: with debate disabled the system no longer
\emph{self-reports} a hallucination rate. That rate must therefore be
measured by the same external judge used for every other condition. Reporting
a self-computed zero for an unverified configuration would be a measurement
artifact of exactly the kind this paper argues against.

\subsection{RQ4: Is the Internal Verifier Trustworthy?}

This question must be answered before any other number is believed, because
the verifier and the generator share a model family. Table~\ref{tab:agreement}
cross-tabulates the internal verdict against human raters and against external
model judges.

\begin{table}[!t]
\centering
\caption{Agreement on hallucination judgments (per claim).}
\label{tab:agreement}
\begin{tabular}{@{}lc@{}}
\toprule
\textbf{Comparison} & \textbf{Agreement} \\
\midrule
Internal vs.\ human  & XX.X\% ($\kappa$ = 0.XX) \\
External vs.\ human  & XX.X\% ($\kappa$ = 0.XX) \\
Internal vs.\ external & XX.X\% ($\kappa$ = 0.XX) \\
Human--human (rater pair) & XX.X\% ($\kappa$ = 0.XX) \\
\bottomrule
\end{tabular}
\end{table}

\todo{Interpretation to be written once numbers exist. If internal--human
agreement is materially lower than human--human agreement, the internal
verifier over- or under-reports, and the headline hallucination rate must be
restated using the human-anchored estimate. Do not suppress this case.}

\subsection{RQ5: Human Evaluation}

\begin{table}[!t]
\centering
\caption{Human ratings (5-point Likert, mean $\pm$ SD) on \todo{30} bundles.}
\label{tab:human}
\begin{tabular}{@{}lcc@{}}
\toprule
\textbf{Dimension} & \textbf{\sysname{}} & \textbf{Best baseline} \\
\midrule
Factual correctness      & X.XX $\pm$ X.XX & X.XX $\pm$ X.XX \\
Difficulty fit           & X.XX $\pm$ X.XX & X.XX $\pm$ X.XX \\
Instructional usefulness & X.XX $\pm$ X.XX & X.XX $\pm$ X.XX \\
\bottomrule
\end{tabular}
\end{table}

\subsection{Cost and Latency}

\begin{table}[!t]
\centering
\caption{Computational cost per case (mean).}
\label{tab:cost}
\begin{tabular}{@{}lc@{}}
\toprule
\textbf{Metric} & \textbf{Value} \\
\midrule
Wall-clock per case & \todo{XX.X s} \\
LLM calls per case  & \todo{XX} \\
Tokens per case     & \todo{XX,XXX} \\
Retry rate          & \todo{XX.X\%} \\
\bottomrule
\end{tabular}
\end{table}

The retry rate is worth noting: the bounded retry policy of
Algorithm~\ref{alg:verify} is designed to trigger rarely, so verification
overhead should stay far below a naive always-refine-twice design.
\todo{Confirm against measured retry rate.}

%=====================================================================
%  6. Discussion
%=====================================================================
\section{Discussion}
\label{sec:discussion}

\subsection{What the Results Mean}

\emph{On orchestration (RQ1).} The comparison between B2 and the full system
isolates the value of orchestration itself: both see the same profile, but only
the latter converts it into an explicit, auditable difficulty target before any
text is generated. \todo{State the observed gap.} The mechanism is unlikely to
be raw model capability---both use the same backend---and more likely to be the
removal of a single-prompt bottleneck, in which one prompt must simultaneously
perform learner interpretation and content production.

\emph{On retrieval (RQ2).} Removing retrieval degrades coverage more than it
degrades any other metric, which is consistent with the view that retrieval
supplies \emph{content} while later stages supply \emph{scrutiny}. The high
retrieval support measured in Section~\ref{ssec:materials} confirms that
residual coverage gaps are a generation failure, not an evidence failure.

\emph{On debate versus self-refinement (RQ3).} B5 is the informative
condition. Self-refinement asks the same model to critique itself; the debate
stage instead requires an explicit claim-level verdict with evidence pointers
and forces cross-agent contradiction resolution. \todo{State the observed
difference.} If debate dominates, the operative factor is likely the
\emph{adversarial framing plus structured output requirement} rather than
simply spending more tokens.

\emph{On the internal verifier (RQ4).} We treat this as a first-class result
rather than a caveat, because the credibility of every other number depends on
it. A system that judges its own output with the model that produced it is
structurally at risk of both false reassurance and false alarms. Reporting
internal, external, and human verdicts side by side lets readers discount our
headline figure appropriately.

\subsection{Implications for Practice}

Three properties of the design transfer to real deployments. First, making the
difficulty rule explicit and deterministic means an instructor can inspect and
override it---a requirement in institutional settings where opaque adaptation
is unacceptable. Second, bounding the retry loop keeps cost predictable, which
matters more than marginal quality gains in production. Third, generating the
six modalities from a shared state means a revision propagates: correcting a
misstated concept in the document also corrects the quiz item and video
script that depend on it.

\subsection{Threats to Validity}
\label{ssec:threats}

\emph{Construct validity.} Knowledge-point coverage is measured by keyword
matching, which can overestimate coverage when a term appears incidentally and
underestimate it when a concept is paraphrased. We mitigate by matching on
multi-token keywords extracted from each knowledge point, but the metric
remains a proxy.

\emph{Internal validity.} The most important limitation is that
\textbf{we report no learner-outcome study}. All results concern the
properties of generated artifacts and their fit to profiles, not measured
learning gain. Claims about improved learning would require a controlled
pre-test/post-test study with human participants, which is outside the scope of
the present system paper and is our immediate next step. A second internal
threat is that profiles and topics were constructed by the authors; although
the difficulty rule is public, the topic set may favor domains where our
knowledge base is dense.

\emph{External validity.} All experiments use a single LLM backend, a single
language of instruction, and one vertical domain family. Nothing here
establishes that the orchestration gains transfer to a different model family
or to non-technical subjects.

\emph{Conclusion validity.} Configurations are compared on identical cases with
paired tests, but the number of cases bounds the detectable effect size;
non-significant differences in Table~\ref{tab:main} should be read as
inconclusive rather than as evidence of equivalence.

\subsection{Ethical Considerations}

Learner profiles contain educational and employment information; the system
stores pseudonymized profile fields and separates identity from behavioral
logs. More substantively, AI-generated instructional material in engineering
domains carries a duty of care: an undetected error can propagate into a
learner's practice. This is precisely why the debate stage is positioned
before delivery rather than after, and why we recommend that deployment in
high-stakes settings retain human instructor sign-off on verified-flagged
content.

%=====================================================================
%  7. Conclusion
%=====================================================================
\section{Conclusion}
\label{sec:conclusion}

We presented \sysname{}, a multi-agent framework for personalized domain
knowledge generation in vertical skill training. Its distinctive feature is
structural rather than algorithmic: verification is placed between generation
and delivery, the difficulty rule is explicit and independently recomputable,
and adaptation is expressed as a published decision table rather than left
implicit in model weights. Across \todo{61} end-to-end cases spanning
\todo{14} topics and \todo{12} learner profiles, the full configuration
outperforms single-pass and ablated alternatives on adaptation accuracy,
coverage, and hallucination rate, with the debate stage and retrieval
grounding contributing the largest measurable shares.

We are deliberately explicit about what this paper does \emph{not} establish:
we demonstrate artifact-level quality and profile fit, not improved learning
outcomes. Closing that gap requires a controlled human study, which is the
central item on our agenda, together with replication on a second model family
and a longitudinal deployment that tracks profile evolution.

More broadly, the results suggest that for educational generation, the
bottleneck is no longer whether a model can produce plausible material, but
whether a system can demonstrate that the material is correct for \emph{this}
learner. Architectures that make verification and adaptation inspectable
appear better suited to that requirement than architectures that pursue
quality purely through scale.


\appendices
%=====================================================================
%  Appendix A — Reproducibility: agent contracts (abbreviated)
%=====================================================================
\section{Agent Contracts}
\label{sec:appendix}

Each agent is defined by a system prompt that fixes its role, the schema of
its output, and its temperature. Contracts are abbreviated here; the full
prompts ship with the released code.

\begin{table}[!h]
\centering
\caption{Agent contracts and decoding temperature.}
\label{tab:contracts}
\begin{tabular}{@{}p{1.9cm}p{3.1cm}p{1.4cm}@{}}
\toprule
\textbf{Agent} & \textbf{Required output schema} & \textbf{Temp.} \\
\midrule
Supervisor & \texttt{\{agents[], difficulty, rationale\}} & 0.3 \\
Diagnosis & \texttt{\{knowledge\_gaps[], difficulty\_target, priority\_focus[], adaptation\_strategy\}} & 0.3 \\
MindMap & hierarchical outline & 0.4 \\
Document & prose at target difficulty & 0.5 \\
Reading & cited frontier material & 0.6 \\
Code & runnable, commented code & 0.3 \\
Quiz & graded items with rubric & 0.4 \\
Video & scene-level script & 0.6 \\
Debate & \texttt{\{hallucinations[$\{$type, agent, content, severity, correction$\}$], fact\_check, cross\_validation, debate\_summary\}} & 0.2 \\
Review & \texttt{\{scores, issues[], overall, retry\}} & 0.2 \\
Adapt & \texttt{\{decision, reason, actions[], learning\_feedback\}} & 0.4 \\
\bottomrule
\end{tabular}
\end{table}

The debate agent's schema deserves comment: requiring \texttt{type},
\texttt{severity}, \texttt{correction}, and an evidence pointer for every
finding is what converts ``the model noticed a problem'' into an auditable
record. Without a structured schema, a verifier's objections are not
actionable by downstream components, and the retry decision cannot be
expressed as the deterministic policy of Algorithm~\ref{alg:verify}.


\bibliographystyle{IEEEtran}
\bibliography{refs}

\end{document}
